\vspace{0.5em}\noindent\textbf{Mục tiêu Tuần 4:}

\begin{itemize}
\tightlist
\item
  Huấn luyện mô hình ViMQ để nhận dạng văn bản y tế.
\item
  Tìm cách làm cho mô hình học tốt hơn dù dữ liệu y khoa bị thiếu hụt
  hoặc nhiễu.
\item
  Đánh giá độ chính xác của AI và đưa mô hình vào chạy thực tế.
\end{itemize}

\vspace{0.5em}\noindent\textbf{Các công việc triển khai trong tuần:}

\begingroup
\small
\setlength{\tabcolsep}{3pt}
\renewcommand{\arraystretch}{1.15}
\begin{longtable}[]{@{}
  >{\raggedright\arraybackslash}p{(\columnwidth - 8\tabcolsep) * \real{0.2000}}
  >{\raggedright\arraybackslash}p{(\columnwidth - 8\tabcolsep) * \real{0.2000}}
  >{\raggedright\arraybackslash}p{(\columnwidth - 8\tabcolsep) * \real{0.2000}}
  >{\raggedright\arraybackslash}p{(\columnwidth - 8\tabcolsep) * \real{0.2000}}
  >{\raggedright\arraybackslash}p{(\columnwidth - 8\tabcolsep) * \real{0.2000}}@{}}
\toprule\noalign{}
\begin{minipage}[b]{\linewidth}\raggedright
Thứ
\end{minipage} & \begin{minipage}[b]{\linewidth}\raggedright
Công việc
\end{minipage} & \begin{minipage}[b]{\linewidth}\raggedright
Ngày bắt đầu
\end{minipage} & \begin{minipage}[b]{\linewidth}\raggedright
Ngày hoàn thành
\end{minipage} & \begin{minipage}[b]{\linewidth}\raggedright
Nguồn tài liệu
\end{minipage} \\
\midrule\noalign{}
\endhead
\bottomrule\noalign{}
\endlastfoot
\textbf{2} & Cấu hình các thông số cơ bản (Optimizer, hàm Loss) để bắt
đầu quá trình huấn luyện AI. & 22/06/2026 & 22/06/2026 &
\href{https://github.com/tadeephuy/ViMQ}{ViMQ Repository} \\
\textbf{3} & Xây dựng quy trình tự động huấn luyện lặp lại nhiều lần để
AI học được nhiều từ vựng mới hơn. & 23/06/2026 & 23/06/2026 &
\href{https://github.com/tadeephuy/ViMQ}{ViMQ Repository} \\
\textbf{4} & Áp dụng các thuật toán tạo nhiễu ngẫu nhiên giúp mô hình
trở nên mạnh mẽ và ít bị dự đoán sai hơn. & 24/06/2026 & 25/06/2026 &
\href{https://github.com/tadeephuy/ViMQ}{ViMQ Repository} \\
\textbf{5} & Viết các tập lệnh để chấm điểm và kiểm tra xem mô hình AI
đoán đúng được bao nhiêu phần trăm. & 26/06/2026 & 26/06/2026 &
\href{https://github.com/tadeephuy/ViMQ}{ViMQ Repository} \\
\textbf{6} & Đóng gói mô hình AI vừa huấn luyện vào hệ thống chatbot để
phân tích câu hỏi bệnh nhân theo thời gian thực. & 27/06/2026 &
28/06/2026 & \href{https://github.com/tadeephuy/ViMQ}{ViMQ
Repository} \\
\end{longtable}
\endgroup

\vspace{0.5em}\noindent\textbf{Kết quả đạt được:}

\begin{itemize}
\tightlist
\item
  Huấn luyện và tối ưu hóa thành công mô hình NER tiếng Việt chuyên khoa
  y tế ViMQ với chiến lược Self-training tiên tiến.
\item
  Đạt chỉ số F1-Score ấn tượng (88.4\%), vượt trội so với các phương
  pháp trích xuất thực thể truyền thống.
\item
  Tích hợp thành công ViMQ vào lõi AI Backend của Smart Healthcare
  Platform, hoàn thiện bước chuẩn bị ngữ cảnh chuyên sâu cho RAG
  Pipeline.
\end{itemize}
